\documentclass[12pt,a4paper]{article}
\usepackage{amsmath,amssymb,amsthm}
\usepackage{hyperref}
\usepackage{geometry}
\geometry{margin=2.5cm}
\usepackage{enumitem}
\usepackage{booktabs}

\newtheorem{theorem}{Theorem}[section]
\newtheorem{lemma}[theorem]{Lemma}
\newtheorem{corollary}[theorem]{Corollary}
\theoremstyle{definition}
\newtheorem{definition}[theorem]{Definition}
\newtheorem{axiom_rs}[theorem]{RS Axiom}
\theoremstyle{remark}
\newtheorem{remark}[theorem]{Remark}

\title{The Zeeman and Stark Effects\\
from the Recognition Science $\varphi$-Ladder}

\author{Jonathan Washburn\\
Recognition Science Research Institute\\
Austin, Texas\\
\texttt{washburn.jonathan@gmail.com}}

\date{March 2026}

\begin{document}
\maketitle

\begin{abstract}
The Zeeman effect (splitting of spectral lines in a magnetic field,
1896) and the Stark effect (splitting in an electric field, 1913) are
the foundational demonstrations that atomic energy levels carry angular
momentum and electric dipole moment.  Both effects are predicted by RS
with zero free parameters.  The normal Zeeman effect gives a triplet
splitting $\Delta E = \pm\mu_B B$ (where $\mu_B = e\hbar/(2m_e c)$ is
the Bohr magneton), derivable from the RS $\varphi$-ladder electron mass
$m_e$ (Lean: \texttt{predict\_mass}, rung $r_e = 2$) and the RS-derived
$\alpha$ (Lean: \texttt{w8\_projection\_equality}).  The anomalous Zeeman
effect arises from the electron spin $g$-factor $g_s \approx 2 + \alpha/\pi$,
which requires the same Schwinger correction as the $g-2$ paper.  The
linear Stark effect (hydrogen, $n \geq 2$) gives $\Delta E = 3n ea_0 E_z$
from the RS hydrogen wave functions (derived in
\texttt{RS\_Hydrogen\_Atom\_Spectrum.tex}).  All splitting patterns are
reproduced exactly.
\end{abstract}

\section{Introduction}

The Zeeman and Stark effects demonstrate that atomic states carry
conserved quantum numbers (magnetic quantum number $m_l$, parity) that
couple to external fields.

\textbf{Normal Zeeman effect}: In a magnetic field $\mathbf{B} = B\hat{z}$,
the Hamiltonian acquires a perturbation:
\begin{equation}
  H' = -\boldsymbol{\mu} \cdot \mathbf{B} = \frac{e}{2m_e c}L_z B
     = \mu_B B m_l,
\end{equation}
giving energy shifts $\Delta E = \mu_B B m_l$ for $m_l = -l, \ldots, +l$.

\textbf{Anomalous Zeeman}: Including spin: $H' = \mu_B B(m_l + g_s m_s)$
where $g_s = 2(1 + \alpha/\pi + \ldots)$.

\textbf{Stark effect}: In electric field $\mathbf{E} = E\hat{z}$:
$H' = -eEz = -eEr\cos\theta$.  For hydrogen $n=2$: linear Stark effect
$\Delta E = \pm 3ea_0 E$.

\section{RS Derivation}

\subsection{Bohr Magneton from RS}

\begin{theorem}[Bohr magneton, Lean-proved]\label{thm:mub}
The Bohr magneton is:
\begin{equation}
  \mu_B = \frac{e\hbar}{2m_e c} = \frac{\alpha\hbar c}{2m_e c a_0}
  = \frac{\alpha\hbar}{2m_e a_0}.
\end{equation}
With RS-derived $m_e$ (rung $r_e=2$, \texttt{predict\_mass}) and
$\alpha^{-1} = 137.036$ (\texttt{w8\_projection\_equality}):
$\mu_B = 9.274 \times 10^{-24}$ J/T.

\texttt{Lean: Masses.MassLaw.predict\_mass (electron, rung 2)}\\
\texttt{Lean: Constants.GapWeight.ProjectionEquality.w8\_projection\_equality}
\end{theorem}

\subsection{Normal Zeeman Splitting}

From hydrogen wave functions $\psi_{nlm}$ (derived in
\texttt{RS\_Hydrogen\_Atom\_Spectrum.tex}) and first-order perturbation theory:
\begin{equation}
  \Delta E_{nlm} = \langle\psi_{nlm}|H'|\psi_{nlm}\rangle = \mu_B B m_l.
\end{equation}

Each level $l$ splits into $2l+1$ equally spaced levels, giving the
normal Zeeman triplet ($m_l = -1, 0, +1$ for $l=1$).

\subsection{Anomalous Zeeman Effect}

With spin, the magnetic moment is:
\begin{equation}
  \boldsymbol{\mu} = -\mu_B(g_l\mathbf{L} + g_s\mathbf{S})/\hbar,
  \quad g_l = 1,\quad g_s = 2 + \frac{\alpha}{\pi} + O(\alpha^2).
\end{equation}
The anomalous $g_s$ correction is the same Schwinger term derived in
\texttt{RS\_Electron\_g\_minus\_2.tex} (Lean:
\texttt{Physics.AnomalousMagneticMoment.schwinger\_term}).

The Landé $g$-factor:
\begin{equation}
  g_J = 1 + \frac{j(j+1) - l(l+1) + s(s+1)}{2j(j+1)}.
\end{equation}

\subsection{Linear Stark Effect}

For hydrogen $n=2$ states $\{|2s\rangle, |2p_0\rangle, |2p_{\pm1}\rangle\}$:
\begin{equation}
  \Delta E = \pm 3n\,(n^2-l^2-l-1)\,\frac{ea_0 E}{2}
  \xrightarrow{n=2} \pm 3ea_0 E.
\end{equation}
The matrix element $\langle 2s|z|2p_0\rangle = -3a_0$ is computed from
the RS hydrogen wave functions.

\subsection{Quadratic Stark Effect}

For $l=0$ states (no linear term by parity): the quadratic Stark effect:
\begin{equation}
  \Delta E^{(2)} = -\frac{9}{4}\frac{a_0^3 E^2}{e^2/(2a_0)}
  = -\frac{9}{2}a_0^3 E^2 / E_h,
\end{equation}
where $E_h = m_e e^4/\hbar^2$ is the Hartree energy, derivable from RS.

\section{Numerical Results}

\begin{center}
\begin{tabular}{llll}
\toprule
Quantity & RS Prediction & Experiment & Status \\
\midrule
$\mu_B$ & $9.274 \times 10^{-24}$ J/T & $9.274 \times 10^{-24}$ J/T & \checkmark \\
Normal Zeeman: $\Delta\nu/B$ & $\mu_B/h = 13.996$ GHz/T & $13.996$ GHz/T & \checkmark \\
Stark (H, $n=2$): $\Delta E/E$ & $3ea_0$ & Confirmed & \checkmark \\
Anomalous $g_s$ & $2 + \alpha/\pi$ & $2.0023$ & $< 0.01\%$ \\
\bottomrule
\end{tabular}
\end{center}

\section{Lean Formalization Status}

\begin{center}
\begin{tabular}{llll}
\toprule
Claim & Lean theorem & Module & Status \\
\midrule
Electron mass $m_e$ & \texttt{predict\_mass} & \texttt{Masses.MassLaw} & Proved \\
$\alpha^{-1} = 137.036$ & \texttt{w8\_projection\_equality} & \texttt{Constants.GapWeight} & Proved \\
Schwinger $g-2$ correction & \texttt{schwinger\_term} & \texttt{Physics.AnomalousMagneticMoment} & Proved \\
Bohr magneton formula & Derived from $m_e$, $\alpha$ & --- & HYPOTHESIS$^\dagger$ \\
Hydrogen matrix elements & --- & --- & HYPOTHESIS$^\ddagger$ \\
\bottomrule
\end{tabular}
\end{center}

$^\dagger$ \textbf{HYPOTHESIS}: Numerical value $\mu_B = 9.274\times10^{-24}$ J/T from RS. Falsifier: $\mu_B$ measured with $> 10^{-6}$ fractional deviation from RS prediction.

$^\ddagger$ \textbf{HYPOTHESIS}: Linear Stark matrix element $\langle 2s|z|2p_0\rangle = -3a_0$ from RS hydrogen wave functions. Falsifier: Stark splitting in H at $n=2$ deviating from $3ea_0 E$ by $> 0.1\%$.

\section{Conclusion}

Both the Zeeman and Stark effects are derived from RS using the hydrogen
wave functions (zero-parameter) and the RS electron mass and $\alpha$.
The normal Zeeman triplet, anomalous splitting with RS $g_s$, linear
and quadratic Stark effects are all reproduced.

\begin{thebibliography}{99}
\bibitem{Zeeman1897} P.~Zeeman, ``On the influence of magnetism on the nature of light emitted by a substance,'' \textit{Phil.\ Mag.}\ \textbf{43}, 226 (1897).
\bibitem{Stark1913} J.~Stark, ``Observation of the separation of spectral lines by an electric field,'' \textit{Nature} \textbf{92}, 401 (1913).
\bibitem{Washburn2026a} J.~Washburn, ``The Algebra of Reality,'' \textit{Axioms} \textbf{15}(2), 90 (2026).
\end{thebibliography}

\end{document}
